% SUPERSEDED: use preliminaries.tex for the current paper opening.
% This earlier probe-based motivation is retained as an editorial archive.
% ---------------------------------------------------------------------------
% Section 1: motivation and empirical puzzle. ICLR 2027; requires booktabs.
% Citation: section1_references.bib. Editorial audit: section1_review.md.
% Outline: mode-averaging hypothesis -> density evidence -> mean intervention
% -> published regression gap -> question about the training objective.
% ---------------------------------------------------------------------------
\section{Why Does MSE Underperform When One Action Suffices?}
\label{sec:single-mode}

% Opening: state the explanation being tested.
Generative robot policies are often motivated by their ability to represent
multiple valid actions for the same observation. If those actions follow
incompatible strategies, a policy trained with mean squared error (MSE) may
average them into an invalid motion. This explanation predicts that repeated
samples from a successful policy should reveal distinct action strategies.
We test this prediction by holding the observation and instruction fixed and
sampling the policy repeatedly.

% Design and evidence: compare predictive density at the same policy input.
The continuous actions favor a single Gaussian description.
We probe GR00T on GR1 and $\pi_{0.5}$ on LIBERO using their evaluation-time
samplers. For each policy, we screen 3,000 states and retain those with the
most separated samples. We then draw 32 fresh action chunks per retained
state and compare a single Gaussian with a two-component Gaussian mixture
along an independently estimated principal direction. The analysis covers
GR1's arm and waist joints and LIBERO's six arm channels.
Table~\ref{tab:single-mode-density} shows that the median held-out likelihood
favors the single Gaussian in both settings. Even this search for separated
samples supports a simple description of continuous motion: variation around
one local action.

\begin{table}[ht]
\centering
\caption{\textbf{Continuous actions favor a single Gaussian.}
The 1G advantage is the median held-out log-likelihood difference between a
single Gaussian and a two-component mixture, in nats per projected sample.
Positive values favor the single Gaussian. States are selected from a
3,000-state screen to maximize apparent separation; each has 32 fresh samples.}
\label{tab:single-mode-density}
% Source: gaussianity/action_basin_report.json, independent_pc1.state_deltas.
% Average reciprocal split scores within each state, then take the median.
% Original direction-pooled medians were -0.279/-0.266 for 2G minus 1G.
\begin{tabular}{@{}llrr@{}}
\toprule
Policy / benchmark & Channels & States & 1G advantage \\
\midrule
GR00T / GR1          & Arm $+$ waist & 30 & $+0.331$ \\
$\pi_{0.5}$ / LIBERO & Arm           & 36 & $+0.352$ \\
\bottomrule
\end{tabular}
\end{table}

% Intervention: directly test the consequence relevant to averaging.
Averaging the sampled actions also preserves the local motion.
At 40 WidowX states, we estimate the arm-action mean from 64 samples, restore
the simulator state, and execute that mean over the four control steps used
before replanning. Each comparison holds the gripper commands fixed.
The resulting motion stays inside the empirical 95\% envelope of sampled
paths at all 40 states. Thus, in this direct intervention, averaging produces
an ordinary local motion. A regression policy has a useful action to predict.

% Challenge: state the published performance gap with one reporting metric.
Yet MSE regression can still perform substantially worse than Flow.
Table~12 of \citet{pan2025muchado} reports a clear gap on state-based
manipulation tasks with the same backbone and training data. With Sudeep-DiT,
best-checkpoint success is 12\% for regression versus 40\% for Flow on
Transport-mh, and 52\% versus 86\% on Tool-Hang, averaged over three seeds.
The Transport-mh gap also appears with Chi-Transformer and Chi-UNet.
The difficulty therefore persists even when regression uses the same
architecture as the generative policy.

% Transition: distinguish a useful prediction from learning it reliably.
These observations pose a simple question:
\emph{if one local action suffices, why is MSE an ineffective way to learn it?}
MSE targets the conditional mean, making it a natural objective for a
Gaussian action model. The unresolved issue is how the training loss fits
that prediction across the dataset. We next examine two properties that
shape this process: action distributions with different radii at different
task stages, and large errors from rare or noisy actions. These properties
motivate a loss that adapts to both local scale and heavy tails.
